\section{Conclusion}

A multi-level evaluation of 375 model configurations across five genotoxicity datasets shows that, for the Ames mutagenicity benchmark studied here, source heterogeneity---the mixing of drug-like and industrial compounds with sharply different prevalence and provenance---drives a substantially larger performance gap ($\Delta = 0.390$, scaffold CV $\to$ LDO) than the choice between random and scaffold splitting ($\Delta = 0.046$). Threshold-tuning analysis attributes only part of this gap to prior shift, leaving a residual that no recalibration can fix. Algorithm differences are small in-domain (MCC range = 0.063) but become pronounced under source shift, with Random Forest more robust than gradient boosting. In~vivo micronucleus models show ranking signal but lack reliable classification utility with current 2D descriptors. We recommend multi-level evaluation with source-awareness as a minimum reporting standard for Ames mutagenicity QSAR.
